\chapter{إعادة فحص الأقفال والتزامن المتقدم}
\label{CH:LOCK2}

بعد دراسة جميع مكونات النواة (الذاكرة، المصايد، الجدولة، والمقاطعات، ونظام الملفات)، يقدم هذا الفصل مراجعة شاملة وإعادة تفحص متعمقة لكيفية تفاعل الأقفال وميكانيكيات التزامن المتقدمة لحماية هياكل البيانات في \texttt{xv6}.

\section{أنماط الأقفال والتنسيق}
\label{sec:locking_patterns}

يحتوي المخزن المؤقت للكتل (Buffer Cache) في \texttt{kernel/bio.c} على مستويين من الأقفال:
\begin{enumerate}
  \item \textbf{القفل العام (\texttt{bcache.lock})}: يحمي القائمة المتصلة المزدوجة للمخزن المؤقت (\texttt{bcache.head}).
  \item \textbf{القفل الفردي للكتلة (\texttt{b->lock})}: قفل نوم (\texttt{sleeplock}) مخصص لكل كتلة مؤقتة حية.
\end{enumerate">

\section{أنماط التزامن الشبيهة بالأقفال}
\label{sec:lock_like_patterns}

تتبع العقد الفهرسية (Inodes) نمطًا مشابهًا للمخزن المؤقت:
\begin{itemize}
  \item \textbf{القفل العام جدول العقد (\texttt{icache.lock})}: يحمي القائمة الموزعة للجدول وتخصيص المداخل في الذاكرة.
  \item \textbf{القفل الفردي للعقدة (\texttt{ip->lock})}: قفل نوم يحمي حقول العقدة الفهرسية ومحتويات كتلها المربوطة.
\end{itemize}

\section{كود بدون أقفال والتعليمات الذرية}
\label{sec:no_locks_at_all}

في بعض الحالات النادرة مثل السجل العتادي \texttt{tp} الخاص بمعرف وحدة المعالجة، يمكن الاستغناء عن الأقفال الصريحة بشريط تعيين ذرية محكومة عتاديًا.

\section{التوازية والتوسع الأفقية}
\label{sec:parallelism}

تستخدم النواة القفل المخصص لكل عملية (\texttt{p->lock}) لحماية حالة العملية (\texttt{p->state}) ومستقبل القناة (\texttt{p->chan}).

\section{تمارين}
\label{sec:exercises_lock2}

\begin{enumerate}
  \item اشرح لماذا يُفضل استخدام أقفال النوم (Sleep-locks) لحماية كتل القرص بدلاً من أقفال التناوب (Spinlocks)؟
  \item تتبع ترتيب الأقفال المحوزة أثناء تنفيذ استدعاء النظام \texttt{sys\_link()} لمنع حدوث الجمود.
\end{enumerate}
